\section{Câu hỏi nghiên cứu và giả thuyết}\label{sec:questions}

Phần hồi cứu trả lời RQ1 và RQ2. Phần thực nghiệm trả lời RQ3 đến RQ6.

\begin{itemize}
\item \textbf{RQ1.} AI đã tham gia từng khâu của vòng đời ViPharmind ở mức nào, và để lại dấu vết gì
trong kho mã?
\item \textbf{RQ2.} Mã do AI viết phải sửa lại ở đâu và bao nhiêu?
\item \textbf{RQ3.} Trên cùng một nhiệm vụ, làm theo đặc tả (viết thiết kế và kế hoạch trước) có cho kết
quả tốt hơn giao việc trực tiếp (vibe coding) không, và tốn thêm bao nhiêu chi phí?
\item \textbf{RQ4.} Kết luận của RQ3 có giữ nguyên khi đổi sang trợ lý lập trình và mô hình của hãng
khác không?
\item \textbf{RQ5.} Lợi ích của làm theo đặc tả thay đổi thế nào khi dùng mô hình mạnh hơn (Claude Opus
5.5 so với Claude Sonnet 5.5, cùng trong Claude Code)?
\item \textbf{RQ6 (khám phá).} Khoảng cách giữa hai cách làm thay đổi thế nào khi nhiệm vụ lớn cỡ một
tính năng (nhiều commit, khoảng 750 đến 1.500 dòng) thay vì một commit?
\end{itemize}

Bốn giả thuyết được ghi trước khi chạy (mục~\ref{sec:method}, phần chốt trước khi chạy):
\begin{itemize}
\item \textbf{H1.} Làm theo đặc tả có tỉ lệ đạt nhiệm vụ cao hơn vibe coding.
\item \textbf{H2.} Làm theo đặc tả tốn nhiều token và thời gian hơn.
\item \textbf{H3.} Làm theo đặc tả có ít lỗi chất lượng tĩnh và ít hồi quy hơn, và chỉ số bảo trì không
kém hơn.
\item \textbf{H4.} Khoảng cách về tỉ lệ đạt giữa hai cách làm ở Sonnet lớn hơn ở Opus, tức làm theo đặc
tả giúp mô hình yếu hơn nhiều hơn.
\end{itemize}

RQ6 là câu hỏi khám phá. Tầng tính năng chỉ có \NumFeatureTasks{} nhiệm vụ, nên phần trả lời chỉ mô tả
kèm khoảng tin cậy, không kiểm định giả thuyết.
